\section{Multimodal Language Models and LLM Applications}
\label{sec:mllm}

Large language models (LLMs) show strong reasoning abilities in dialogue,
step-by-step argumentation, math problems, and code writing. They learn all of
this from text alone. A robot that should ``bring me the rice chips from the
drawer'' needs more: it must know where the drawer is, whether it is open, and
which object in the camera image is the bag of chips. This is the
\emph{grounding problem}. LLM representations must be connected to what the
system perceives (images, sensor states), to what it can do (actions in the
physical world), and to facts about the concrete scene, such as its geometric
configuration. Text-only training never supplies this link between words and
percepts.

This section treats two answers to the grounding problem. The first is
architectural: PaLM-E injects continuous sensor observations directly into the
embedding space of a pretrained LLM and trains the combination end-to-end, so
that one model plans robot actions, answers visual questions, and writes
captions. The second is practical: an application program calls an existing
vision-capable LLM through an API and wraps it in a small workflow graph that
decides what to do with each image. Both build on the decoder-only transformer
of Section~\ref{sec:transformers} and on the Vision Transformer (ViT) of
Equation~\ref{eq:vit-patches}.

\subsection{Decoder-Only Language Models and Prompts}

A decoder-only LLM is a generative model of text. Text is represented as a
token sequence $w_{1:L}=(w_1,\dots,w_L)$, and the model assigns it a
probability by the autoregressive factorization of
Equation~\ref{eq:transformer-autoregressive},
\begin{equation}
  p(w_{1:L}) = \prod_{l=1}^{L} p_{\mathrm{LM}}(w_l \mid w_{1:l-1}) .
  \label{eq:mllm-lm}
\end{equation}
\noindent
Here $p_{\mathrm{LM}}$ is the transformer with causal masking: the output at
position $l$ is a softmax over the vocabulary and predicts the next token.
Before the first layer, every token passes through the embedding table,
written here as a function $\gamma:\mathcal W\to\mathbb R^k$ that maps a
vocabulary entry $w$ to a vector $x=\gamma(w)$ in a $k$-dimensional
embedding space.

A \emph{prefix-decoder} LLM splits the sequence into a given prefix
$w_{1:n}$ and a continuation that the model predicts:
\begin{equation}
  p(w_{n+1:L}\mid w_{1:n})
    = \prod_{l=n+1}^{L} p_{\mathrm{LM}}(w_l \mid w_{1:l-1}) .
  \label{eq:mllm-prefix}
\end{equation}
\noindent
The prefix, often called the \emph{prompt}, provides the context. It may
describe the task (``Q: How can the yellow block be moved onto the blue plate?
A:''), contain earlier dialogue, or show a few solved examples. The LLM
continues the prefix, and that continuation is the answer. Three prompt
regimes follow from this view and appear repeatedly below:
\begin{itemize}
  \item \emph{Zero-shot}: the prefix contains only the task, no solved
    examples. The model must rely on what it learned in pretraining.
  \item \emph{Few-shot} (in-context learning): the prefix contains a few
    input--output examples, and the model continues the pattern. The weights
    do not change; only the input does.
  \item \emph{Chain of thought} (CoT): the prefix asks for, or demonstrates,
    intermediate reasoning steps before the final answer (``Let's think step
    by step''). Spelling out steps turns one hard prediction into several
    easier next-token predictions, each conditioned on the previous steps.
\end{itemize}

\noindent
All three change behavior at inference time without training. This makes a
pretrained LLM a general interface: whoever can express a task as a prefix
can use it. The rest of the section asks what happens when parts of that
prefix are not words but images and sensor readings.

\subsection{PaLM-E: Injecting Observations into the Language Embedding Space}

An \emph{embodied language model} incorporates continuous inputs from the
real-world sensor and actuator modalities of an agent, such as a robot, into
a language model. PaLM-E (``Pathways Language Model, Embodied'') builds it
from PaLM, a large pretrained decoder-only LLM. The central idea is simple.
The LLM already processes a sequence of $k$-dimensional vectors. If an image
or a state estimate is encoded into vectors of the same dimension $k$, it can
be placed into the same sequence as the word embeddings and processed by the
self-attention layers \emph{in the same way as text}. Nothing in the
transformer needs to change.

The input is a \emph{multimodal sentence}: text with interleaved
placeholders for observations, for example
\begin{center}
  \texttt{Q: What happened between <img\_1> and <img\_2>?}
\end{center}
\noindent
Formally, an encoder $\phi_j$ maps observation $O_j$ (an image, a state
vector, a 3D scene) to a sequence of $q_j$ vectors,
$\phi_j(O_j)\in\mathbb R^{q_j\times k}$. The input vector at position $i$ is
then
\begin{equation}
  x_i =
  \begin{cases}
    \gamma(w_i) & \text{if position } i \text{ is a text token},\\
    \phi_j(O_j)_i & \text{if position } i \text{ belongs to observation } O_j .
  \end{cases}
  \label{eq:mllm-multimodal-sentence}
\end{equation}
\noindent
The placeholder \texttt{<img\_1>} is thus replaced by $q_1$ vectors, not by
one token. Observation vectors can sit anywhere between the text tokens, so
the text can refer to an image before or after it and several images can be
compared. Figure~\ref{fig:palme-architecture} shows the complete model.

\begin{figure}[ht]
  \centering
  \begin{tikzpicture}[
    box/.style={draw, rounded corners=2pt, minimum height=0.8cm,
      align=center, font=\scriptsize},
    tok/.style={draw, minimum width=0.42cm, minimum height=0.42cm,
      inner sep=0pt},
    arr/.style={-{Latex[length=2mm]}, thick},
    lbl/.style={font=\scriptsize, align=center}
  ]
    % inputs
    \node[box, fill=blue!10, minimum width=1.6cm] (img) at (0,0)
      {image $I$};
    \node[box, fill=blue!10, minimum width=1.6cm] (vit) at (0,1.35)
      {ViT $+$ $\psi$};
    \node[box, fill=yellow!20, minimum width=2.6cm] (txt) at (4.9,0)
      {text: ``Q: How to grasp\\the white object? A:''};
    \node[box, fill=yellow!20, minimum width=2.6cm] (emb) at (4.9,1.35)
      {token embedder $\gamma$};
    \draw[arr] (img) -- (vit);
    \draw[arr] (txt) -- (emb);
    % token row
    \foreach \i in {0,...,3}
      \node[tok, fill=blue!25] (v\i) at (-0.63+\i*0.42,2.6) {};
    \foreach \i in {0,...,6}
      \node[tok, fill=yellow!45] (t\i) at (3.64+\i*0.42,2.6) {};
    \node[lbl] at (1.5,2.6) {$\cdots$};
    \draw[arr] (vit) -- (0,2.35);
    \draw[arr] (emb) -- (4.9,2.35);
    \draw[decorate, decoration={brace, amplitude=4pt}]
      (-0.85,2.9) -- (5.3,2.9)
      node[midway, above=5pt, lbl] {multimodal sentence
      $x_1,\dots,x_n$ (prefix), all in $\mathbb R^k$};
    % LLM
    \node[box, fill=green!12, minimum width=7.2cm, minimum height=0.9cm]
      (llm) at (2.2,4.35) {decoder-only LLM (PaLM), causal self-attention};
    \draw[arr] (2.2,3.45) -- (llm.south);
    % output
    \node[box, fill=orange!15, minimum width=5.6cm] (out) at (2.2,5.75)
      {``First, grasp the orange object and place it \dots''};
    \draw[arr] (llm) -- (out);
    \node[lbl, above=2pt of out] {output text: answer or plan};
    % robot loop
    \node[box, fill=gray!12, minimum width=1.9cm] (pol) at (8.9,5.75)
      {low-level\\policy};
    \node[box, fill=gray!12, minimum width=1.9cm] (rob) at (8.9,1.35)
      {robot acts,\\camera observes};
    \draw[arr] (out.east) ++(0,0) -- (pol.west);
    \draw[arr] (pol) -- node[right, lbl] {motor\\commands} (rob);
    \draw[arr, dashed] (rob.south) |- node[pos=0.75, below, lbl]
      {new image (closed loop)} (img.east |- 0,-0.9) -- (img.south);
  \end{tikzpicture}
  \caption{PaLM-E. An encoder turns each observation into a few vectors of
  the LLM's embedding dimension $k$. They are interleaved with the word
  embeddings to a multimodal prefix, and the pretrained LLM continues it
  with text. For robot tasks, the text is a plan step that a low-level policy
  executes; the next camera image starts the next planning step.}
  \label{fig:palme-architecture}
\end{figure}

\noindent
The output is always text. For visual question answering (VQA) and
captioning, the text is the answer or description. For robotics, the text is
a sequence of high-level decisions (``pick the green block,'' ``push it to
the corner''). PaLM-E does not output joint angles itself. A separate
\emph{low-level policy}, a smaller learned controller that understands such
short language commands, translates each step into motor commands. After the
step, the robot takes a new image, and PaLM-E is queried again with it. This
\emph{closed-loop} planning lets the model react to failures and
disturbances: if a person moves the chips back into the drawer, the next
image shows it, and the plan adapts.

\paragraph{Worked example: sequence length.}
An image of $224\times224$ pixels encoded by a ViT with patch size $P=14$
gives $(224/14)^2=16^2=256$ patch features
(Equation~\ref{eq:vit-patches}). A learned projection $\psi$ maps each
feature from the ViT width to the LLM width $k$. The image therefore
occupies $q=256$ positions of the prefix. The question ``What happened
between <img\_1> and <img\_2>?'' with about $10$ text tokens becomes a
prefix of roughly $2\cdot256+10=522$ vectors. Images, not words, dominate the
sequence length and thus the quadratic attention cost. This is one reason to
prefer compact encoders, such as the object-centric ones below, that describe
a scene with few vectors.

\subsection{Encoders for Different Sensor Modalities}

Equation~\ref{eq:mllm-multimodal-sentence} leaves the encoder $\phi$ open.
Its job is to map an observation into vectors the LLM can use, and different
modalities call for different encoders. PaLM-E compares several
(Table~\ref{tab:palme-encoders}).

\begin{itemize}
  \item \emph{State estimation vectors.} If the robot or a perception system
    already provides a state $s\in\mathbb R^S$, e.g.\ the poses of all objects,
    a multilayer perceptron maps it into the embedding space,
    $\phi_{\mathrm{state}}(s)\in\mathbb R^{k}$. This is the simplest input but
    requires a working state estimator and fixes which quantities are known.
  \item \emph{Vision Transformer.} A pretrained ViT
    $\tilde\phi_{\mathrm{ViT}}$ maps the image $I$ to $m$ token features. Its
    feature width generally differs from $k$, so a learned affine map $\psi$
    projects each feature:
    $\phi_{\mathrm{ViT}}(I)=\psi\bigl(\tilde\phi_{\mathrm{ViT}}(I)\bigr)
    \in\mathbb R^{m\times k}$. This works on raw pixels and needs no state
    estimator. Its tokens, however, describe image patches, not objects.
  \item \emph{Object-centric representations.} Planning statements are about
    objects (``put the red block on the plate''). Object-centric encoders
    separate the visual input into distinct objects before injecting it into
    the LLM. Conceptually this requires instance masks $M_j$: the object token
    is $\phi_{\mathrm{ViT}}(I\circ M_j)$, the ViT applied to the masked image.
    Ground-truth masks are available in simulation but not in the real world.
  \item \emph{Object Scene Representation Transformer (OSRT).} OSRT learns a
    neural 3D scene representation from several views without supervision.
    It discovers objects by itself and represents the scene by a set of
    \emph{object slots} $o_j$. Each slot is projected by an MLP to one or
    more tokens. This gives object-centric, 3D-aware inputs without
    ground-truth masks.
\end{itemize}

\noindent
Object-centric encoders enable \emph{entity referrals}. The prefix can name
each object with its own tokens, e.g.\ ``Object 1 is <obj\_1>. Object 2 is
<obj\_2>. \dots'' The model can then answer ``Push obj\_2 to the left'' and
the low-level policy knows exactly which object is meant. A plain ViT
encoding offers no such handle; the plan must describe objects by color or
position, which is ambiguous when several objects look alike.

\begin{table}[ht]
  \centering
  \begin{tabular}{@{}p{2.6cm}p{3.1cm}p{3.6cm}p{4.2cm}@{}}
    \toprule
    Encoder & Input & Tokens produced & Main strength / limitation \\
    \midrule
    State MLP & state vector $s$ & one vector per state & compact and exact;
      needs a state estimator \\
    ViT $+\psi$ & raw image & one per patch (e.g.\ 256) & general, pretrained
      at scale; not object-based \\
    Object-centric ViT & image and masks $M_j$ & tokens per object & supports
      entity referrals; needs masks \\
    OSRT & multi-view images & tokens per object slot & 3D-aware, objects
      learned without masks; most data-efficient \\
    \bottomrule
  \end{tabular}
  \caption{Input encoders compared in PaLM-E. All map an observation to
  vectors of the LLM embedding dimension $k$; they differ in what structure
  they impose on the scene.}
  \label{tab:palme-encoders}
\end{table}

\subsection{Training PaLM-E and Its Design Choices}

PaLM-E is trained end-to-end on a dataset of multimodal examples. Example $i$
consists of $u_i$ continuous observations $O^i_{1:u_i}$, a token sequence
$w^i_{1:L_i}$ with placeholders for them, and an index $n_i$ that marks where
the prefix ends. The loss is the ordinary next-token cross-entropy, but only
on the continuation:
\begin{equation}
  \mathcal L = -\sum_i \sum_{l=n_i+1}^{L_i}
    \log p_{\mathrm{LM}}\bigl(w^i_l \mid x^i_{1:l-1}\bigr),
  \label{eq:mllm-loss}
\end{equation}
\noindent
where $x^i$ is the multimodal sentence of
Equation~\ref{eq:mllm-multimodal-sentence}. The prefix, containing the
images and the question, is only conditioned on; the model is never trained
to predict it. Observation vectors are not tokens from the vocabulary, so
there is no target to predict at those positions anyway. Gradients flow
through the LLM into the projection $\psi$ and the encoder, which learn to
produce vectors that the LLM can ``read.''

\paragraph{Worked example: which tokens are trained.}
Take the training example ``<img> Q: What color is the cube? A: red.'' with
$256$ image vectors, $8$ question tokens, and $2$ answer tokens (``red,''
``.''). Then $n_i=264$, and the loss sums over just the two answer positions.
If the model predicts $p(\text{red}\mid\text{prefix})=0.5$ and
$p(\text{.}\mid\text{prefix, red})=0.9$, the example contributes
$-\log0.5-\log0.9\approx0.69+0.11=0.80$. The $264$ prefix positions still
matter: every answer prediction attends to them, so their encoders receive
gradients through attention.

\paragraph{Frozen or fine-tuned LLM.}
Two training regimes are possible:
\begin{itemize}
  \item \emph{Full fine-tuning}: encoder, projection, and LLM weights are all
    updated. The model can adapt its language layers to the new inputs but
    may lose some of its original abilities.
  \item \emph{Frozen LLM}: only the encoder $\phi$ and projection $\psi$ are
    trained. The encoder must learn to translate observations into vectors
    that the fixed LLM already understands. Because these vectors act like a
    learned, image-dependent prompt, this is called \emph{input-conditioned
    soft prompting}. It keeps the language abilities intact by construction.
\end{itemize}

\noindent
Encoders can themselves be pretrained (a large ViT, a pretrained OSRT) or
trained from scratch. The data question is \emph{sample complexity}: robot
data are expensive to collect, often only a few hundred demonstrations per
task, while images with text are abundant on the web.

\paragraph{Multi-task training and transfer.}
PaLM-E is trained on a \emph{mixture} of tasks: robot planning in several
environments (tabletop task and motion planning, the Language-Table pushing
environment, mobile manipulation in a kitchen), general vision--language data
(VQA, captioning), and pure language tasks. A key finding is \emph{positive
transfer}: the model trained on the full mixture solves each robot task
better than a model trained on that task alone. Knowledge from web-scale
vision--language data carries over to robotics, so a robot task needs
considerably fewer of its own examples. Transfer thus yields high data
efficiency for robotics, and the same model supports zero-shot multimodal
CoT, few-shot prompting, OCR-free math reasoning from images, and
multi-image reasoning.

\paragraph{Scaling and catastrophic forgetting.}
PaLM-E was built at several sizes by combining a PaLM LLM with a ViT, e.g.\
an 8B-parameter LLM with a 4B ViT (PaLM-E-12B) up to a 540B LLM with a 22B
ViT (PaLM-E-562B). When the LLM is fine-tuned on multimodal data, it can
\emph{catastrophically forget} its pure-language abilities: new gradients
overwrite what the old weights encoded. PaLM-E observed that this forgetting
shrinks strongly with scale. The smallest model lost most of its
language-generation performance, the largest only a few percent. Larger
models have enough capacity to add the new modality without giving up the
old skills. Freezing the LLM avoids forgetting entirely, at the price of less
adaptation.

\paragraph{Emergent capabilities.}
Some abilities appear in the largest model without having been trained
explicitly. With \emph{zero-shot multimodal chain of thought}, the model,
given only the prompt, spells out a step-by-step reasoning narrative that
mixes textual and visual evidence before stating its final answer, e.g.\
``The image shows two bottles of water and one of juice. \dots\ So the answer
is three.'' With \emph{multi-image reasoning}, it compares several images in
one prompt, although training examples mostly contained single images. It
also shows \emph{zero-shot generalization} in robotics to unseen object
pairings or unseen objects. Table~\ref{tab:palme-summary} collects the
research questions and the answers PaLM-E gives.

\begin{table}[ht]
  \centering
  \begin{tabular}{@{}p{5.6cm}p{8.2cm}@{}}
    \toprule
    Question & Finding \\
    \midrule
    Can pixels or states be input to an LLM while keeping its
      generalization? & yes, encoded into the word embedding space and
      interleaved with text \\
    Which input representation works best? & object-centric, 3D-aware (OSRT)
      is most data-efficient; ViT is the general choice for raw images \\
    Frozen or fine-tuned LLM? & frozen LLM with trained encoder already works
      (soft prompting); fine-tuning adapts more but risks forgetting \\
    Single-task or mixture training? & mixture gives positive transfer and
      higher data efficiency on robot tasks \\
    Effect of scale? & less catastrophic forgetting of language; emergent
      multimodal CoT and multi-image reasoning \\
    \bottomrule
  \end{tabular}
  \caption{Research questions of PaLM-E and its answers. One generalist
  model covers robot planning, vision--language, and language tasks.}
  \label{tab:palme-summary}
\end{table}

\noindent
PaLM-E fits the general notion of a foundation model, a model trained on
broad data and adapted to many tasks (Appendix~\ref{app:foundation}). Its
vision encoder is exactly such a pretrained model, and the LLM is another.
PaLM-E shows how two foundation models from different modalities can be
joined through a shared embedding space.

\subsection{Using an LLM Through an API: A Visual Task Agent}

Training a model like PaLM-E requires resources few groups have. In practice,
vision-capable LLMs are available as services: an application sends images
and text over an HTTP API and receives text back. Such a model can already
describe images, read text in them, and answer questions. What remains for
the application developer is to decide \emph{what} to ask, and \emph{what to
do} with the answer. A program that makes these decisions, possibly calling
other tools, is called an \emph{agent}.

Consider a \emph{visual task helper} that receives images of different
kinds and must treat them differently:
\begin{itemize}
  \item a \emph{family image} should get a friendly scene summary;
  \item a \emph{product image} should trigger a secondary tool, e.g.\ a
    search API queried with the detected objects;
  \item a \emph{document or sign} should go to optical character recognition
    (OCR) so that its text is extracted verbatim.
\end{itemize}

\noindent
The flow is: take an image (a local file or a webcam frame); let a vision
model describe it; let the agent decide, based on that description, whether
to summarize the scene, extract text, or call a tool; return the result.
Figure~\ref{fig:vision-agent} shows this flow as a graph.

\begin{figure}[ht]
  \centering
  \begin{tikzpicture}[
    box/.style={draw, rounded corners=2pt, minimum height=0.8cm,
      minimum width=1.9cm, align=center, font=\scriptsize},
    dec/.style={draw, diamond, aspect=1.8, inner sep=1pt, align=center,
      font=\scriptsize, fill=yellow!20},
    arr/.style={-{Latex[length=2mm]}, thick},
    lbl/.style={font=\scriptsize, align=center}
  ]
    \node[box, fill=blue!10] (in) at (0,0) {image\\(file, webcam)};
    \node[box, fill=green!12] (desc) at (2.8,0) {vision LLM:\\describe};
    \node[dec] (route) at (5.7,0) {route};
    \node[box] (sum) at (9.0,1.3) {summarize\\scene};
    \node[box] (ocr) at (9.0,0) {OCR: extract\\text};
    \node[box] (tool) at (9.0,-1.3) {tool call\\(search API)};
    \node[box, fill=orange!15] (out) at (11.9,0) {answer};
    \draw[arr] (in) -- (desc);
    \draw[arr] (desc) -- node[above, lbl] {state} (route);
    \draw[arr] (route.north) |- node[pos=0.75, above, lbl] {family} (sum.west);
    \draw[arr] (route.east) -- node[above, lbl] {document} (ocr.west);
    \draw[arr] (route.south) |- node[pos=0.75, below, lbl] {product}
      (tool.west);
    \draw[arr] (sum.east) -- (out);
    \draw[arr] (ocr.east) -- (out);
    \draw[arr] (tool.east) -- (out);
  \end{tikzpicture}
  \caption{Visual task agent as a workflow graph. Nodes are processing
  steps, which may call an LLM or a tool. The routing node reads the shared
  state, here the image description, and selects the branch.}
  \label{fig:vision-agent}
\end{figure}

\paragraph{LangChain and LangGraph.}
Two open-source frameworks support such programs. \emph{LangChain} provides
building blocks for LLM-powered applications and agents: uniform wrappers
around model providers, message and prompt templates, output parsers, and
tool definitions. \emph{LangGraph} orchestrates LLM and agent workflows as a
\emph{graph}. Each node is a function that reads and updates a shared
\emph{state} (here: the image, its description, the chosen route, the
result). Edges fix the order of nodes; \emph{conditional edges} call a
routing function on the state and branch accordingly. Loops are allowed, so
an agent can also retry or ask a follow-up question. Making the control flow
an explicit graph keeps it inspectable and testable, instead of hiding every
decision inside one long prompt.

\paragraph{A request to a vision LLM.}
A chat API receives a list of \emph{messages}, each with a role
(\texttt{system} for standing instructions, \texttt{user} for the request,
\texttt{assistant} for earlier model replies). A user message can mix text
and images; a local image is usually sent base64-encoded. The core of the
``describe'' node and the graph wiring look like this in Python:

\begin{small}
\begin{verbatim}
from openai import OpenAI
client = OpenAI()               # reads OPENAI_API_KEY from the environment

def describe(state):
    b64 = state["b64"]          # base64-encoded JPEG
    msg = [{"role": "user", "content": [
        {"type": "text", "text": "Describe the image. Is it a family photo,"
                                 " a product, or a document?"},
        {"type": "image_url",
         "image_url": {"url": f"data:image/jpeg;base64,{b64}"}}]}]
    r = client.chat.completions.create(model="gpt-4o-mini", messages=msg)
    return {"description": r.choices[0].message.content}

from langgraph.graph import StateGraph, END
g = StateGraph(State)                       # State: a typed dict
g.add_node("describe", describe)
g.add_node("summarize", summarize); g.add_node("ocr", ocr)
g.add_node("search", search)
g.set_entry_point("describe")
g.add_conditional_edges("describe", route,  # route(state) -> key
    {"family": "summarize", "document": "ocr", "product": "search"})
for n in ("summarize", "ocr", "search"): g.add_edge(n, END)
app = g.compile(); result = app.invoke({"b64": image_b64})
\end{verbatim}
\end{small}

\noindent
Here \texttt{route} can be a plain keyword test on the description or a
second LLM call that answers with one of the three keys. The multimodal
message is the practical counterpart of PaLM-E's multimodal sentence: the
provider's model turns the image into embedding vectors and places them in
the prompt next to the text tokens. The developer only sees text in and text
out.

\paragraph{API keys and running the agent.}
The service authenticates each request with a secret \emph{API key}. It is
created on the provider's developer platform (for billing, a small prepaid
balance suffices) and stored in a \texttt{.env} file in the project folder,
e.g.\ \texttt{OPENAI\_API\_KEY="sk-..."}. The program loads it into the
environment (e.g.\ with \texttt{python-dotenv}); the client library reads it
from there. The key must never appear in code or in a public repository:
anyone who has it can make requests billed to the owner. Add \texttt{.env} to
\texttt{.gitignore} and, where possible, create the key with restricted
permissions. A reproducible Python environment is set up with conda,
\begin{small}
\begin{verbatim}
conda env create -f environment.yml
conda activate vision-agent
python main.py
\end{verbatim}
\end{small}
\noindent
and the agent then runs on local images or the webcam.

\paragraph{Practical limits.}
The API model is a black box that can change between versions, and its
answers are not deterministic. Routing decisions should therefore be
constrained, e.g.\ by asking for exactly one of a fixed set of labels and
falling back to a default otherwise. Every call costs money and time, and
images consume many input tokens (compare the worked example on sequence
length above), so large batches should be sent at reduced resolution when
the task allows. Finally, images sent to an external service leave the local
machine, which matters for private family photos or confidential documents.

\subsection{Trained Integration Versus API Composition}

PaLM-E and the visual task agent answer the grounding problem at different
levels. PaLM-E changes the model: observations become part of the LLM input,
and training teaches the model to reason about them and to plan actions. The
agent keeps the model fixed and changes the program around it: a graph
decides which question to ask and which tool to call.
Table~\ref{tab:mllm-compare} contrasts the two.

\begin{table}[ht]
  \centering
  \begin{tabular}{@{}p{3.2cm}p{5.3cm}p{5.3cm}@{}}
    \toprule
    Aspect & Embodied multimodal LLM (PaLM-E) & LLM via API (visual agent) \\
    \midrule
    Where images enter & encoder $\phi$, vectors in the embedding space &
      image part of a user message; encoding hidden \\
    What is trained & encoder, projection, optionally the LLM & nothing;
      behavior set by prompts and graph \\
    Output & text: answers, captions, robot plans & text, plus results of
      tool calls \\
    Control flow & learned, inside the model (closed-loop replanning) &
      explicit graph with nodes and conditional edges \\
    Requirements & large compute, multimodal and robot data & API key,
      budget, network access \\
    Typical use & research, robotics, generalist models & applications
      built quickly from existing models \\
    \bottomrule
  \end{tabular}
  \caption{Two ways to connect language models with images. Trained
  integration grounds the model itself; API composition grounds an
  application built around a fixed model.}
  \label{tab:mllm-compare}
\end{table}

\noindent
Both views share the same core mechanism. An image is turned into a set of
vectors, those vectors join the text tokens in one sequence, and a
decoder-only transformer continues that sequence with text. Everything
else, whether the text is a VQA answer, a robot plan, or a routing label,
depends on the prompt and on what the text is used for afterwards. Image
\emph{captioning}, one of the basic computer vision tasks, is the simplest
instance: the prefix is an image and ``Describe the image,'' and the
continuation is the caption.
